\section{Conclusion}

This paper examines whether AI adoption is associated with macroeconomic performance in a multi-country panel. Across baseline, robustness, and sensitivity specifications, the estimated association between AI and TFP is positive and statistically stable. By contrast, the estimated relationship with contemporaneous GDP growth is positive but smaller, consistent with delayed macro transmission.

\subsection{Research Questions Revisited}

\textbf{Answer to RQ1:} Qualified yes. Within-country increases in the composite AI/digital index are consistently associated with higher TFP across baseline and robustness specifications in this dataset. However, extended falsification checks indicate this association is driven by the digital-infrastructure component of the index rather than its innovation/patent component, and a reverse-causality check finds that lagged TFP itself positively predicts current AI adoption. The result is therefore best described as a robust association between digital capacity and TFP, with timing that cannot be attributed to AI adoption alone.

\textbf{Answer to RQ2:} Partially. The association with contemporaneous GDP growth is positive but smaller in magnitude and less stable than the TFP result, which is consistent with slower macro transmission -- though, as with RQ1, the reverse-causality evidence means this should not be read as a one-directional timing pattern.

Two contributions are central. Empirically, the positive digital/AI-TFP association persists across alternative fixed-effects structures, outlier handling, lagged regressors, inference choices, and a coverage-restricted sample, but its construct-validity and timing limits are now characterized directly rather than only discussed in principle. Methodologically, all analytical steps are reproducible, and findings -- including the extended falsification checks -- are reported with explicit uncertainty and specification sensitivity.

The interpretation remains deliberately conservative. The analysis identifies robust conditional associations, not definitive causal effects. Remaining concerns include reverse causality (now directly evidenced), construct validity of the AI proxy (now directly evidenced via sub-index decomposition), unobserved time-varying confounders, and a complete-case sample that is concentrated in higher-income, better-governed countries. These limits do not nullify the results, but they define what can credibly be claimed under the current design.

For policy, the findings support a complementary-capabilities perspective: AI adoption is more likely to translate into measurable productivity gains when countries invest in skills, infrastructure, and implementation capacity. For future research, the priority is to pair richer measurement with quasi-experimental designs that can separate causal effects from correlated structural change.

Relative to related studies, these results align with complementarity-based accounts of technology diffusion, while offering a more cautious macro interpretation by showing moderate productivity effects and weaker evidence of immediate aggregate growth acceleration.

In sum, this paper provides a transparent empirical baseline for the AI-productivity debate. Its value lies in demonstrating what appears robust in current data, what remains uncertain, and where stronger evidence is needed.

This positioning is intentional: the manuscript serves as the baseline evidence anchor, while diagnostics-intensive methodological expansion is developed in companion technical work.
